\section{Proof of the Finite-Box Theorem}
\begin{proof}
Write $b_j=2Ba_j-B^2$ and $N_B=\lceil2B/r\rceil+1$. Choose a healthy explanation equal to zero on the prefix. After onset it moves at speed at most $r$ toward the endpoint $\epsilon B$ of the currently observed task, and holds there on arrival. At a task change it starts moving toward the new endpoint. It never has to return to zero at the final time. Interpolation during excluded slots preserves the same rate bound.

For an observed point with target sign $\epsilon$, the deficiency in squared-energy saving relative to $b_j$ is
\begin{equation}
 b_j-(2\epsilon a_jz_j-z_j^2)
 =2a_j(B-\epsilon z_j)-(B^2-z_j^2)\le4Ba_j.
\end{equation}
It is zero once the tracker reaches $\epsilon B$. Following a task change, at most $N_B$ observations can be deficient before the tracker reaches that endpoint; an intervening change restarts the count. Assign each deficient observation to its latest change. Every internal change can be paired with a disjoint complete $k$-slot movement immediately preceding the new observed run, starting at $m$. Its assigned deficiency is at most
\begin{equation}
 4BN_B\{a_m+\eta(k+N_B)\}.
\end{equation}
On an excluded slot the oracle energy exceeds the baseline by $(a_j-B)^2$. The paired move thus supplies at least $k(a_m-B)_+^2$; unpaired excluded slots supply nonnegative excess. A change can therefore have net adverse contribution at most
\begin{equation}
 q_+(a_m)=\pos{4BN_B[a_m+\eta(k+N_B)]-k(a_m-B)_+^2}.
\end{equation}
Because $\eta>0$, $q_+(a_m)$ is nonzero at only finitely many possible start indices. Define the unwhitened constant
\begin{equation}
 C_1=4BN_B(\rho_0+\eta N_B)+\sum_{m\ge1}q_+(\rho_0+\eta m).
\end{equation}
The sum has finitely many nonzero terms. The initial approach from zero contributes the first term, and the tracker stays at its last endpoint. Simultaneously for all calendars,
\begin{equation}
 2\sigma^2D_H^B(\pi)\ge\sum_{j=1}^Hb_j-C_1.\label{eq:finiteLower}
\end{equation}
This is a uniform full-history construction.

For staying, dropping the rate constraints and the nonnegative prefix loss gives
\begin{equation}
 2\sigma^2G_H^B(\mathrm{stay})\ge\sum_{j=1}^H(a_j-B)_+^2.
\end{equation}
Consequently its deficit is at most $\sum b_j/(2\sigma^2)+C_2$, where
\begin{equation}
 2\sigma^2C_2=\sum_{j\ge1:a_j<B}(B-a_j)^2<\infty.
\end{equation}
Conversely, choosing the feasible constant $z=B$ throughout the prefix and subsequent stay experiment gives
\begin{equation}
2\sigma^2D_H^B(\mathrm{stay})\ge\sum_{j=1}^Hb_j-n_0B^2.
\end{equation}
The two inequalities identify its full polynomial baseline up to $O(1)$. Combining the stay upper bound with \eqref{eq:finiteLower}, and using optimality relative to staying, proves both \eqref{eq:finite} and \eqref{eq:regret}.
\end{proof}

An explicit bound in the theorem is
\begin{equation}
 C_B=\frac{C_1+\sum_{j\ge1:a_j<B}(B-a_j)^2}{2\sigma^2}.
 \label{eq:explicit_CB}
\end{equation}
Both terms in the numerator are unwhitened squared-torque quantities. For the fast-ramp example at $\sigma=.2$ this gives about $2.97\times10^5$ nats, compared with a $14.7$-nat advantage in the tested switching family. It is a uniform existence bound, while finite-horizon designs are evaluated by the exact projection.
